\documentclass[10pt,a4paper]{article}

\usepackage[a4paper, top=1.8cm, bottom=1.8cm, left=2cm, right=2cm]{geometry}
\usepackage{amsmath, amssymb}
\usepackage{xcolor}
\usepackage{mdframed}
\usepackage{enumitem}
\usepackage{titlesec}
\usepackage{booktabs}
\usepackage{array}
\usepackage[T1]{fontenc}
\usepackage[utf8]{inputenc}

\usepackage{parskip}


% ── COLORS ──────────────────────────────────────────────────
\definecolor{amber}{HTML}{7A4500}
\definecolor{cyanc}{HTML}{004A7A}
\definecolor{greenc}{HTML}{1A5C1A}
\definecolor{redc}{HTML}{7A1A1A}
\definecolor{purplec}{HTML}{4A1A7A}
\definecolor{greyc}{HTML}{444444}
\definecolor{lightbg}{HTML}{F5F5F5}
\definecolor{bordercol}{HTML}{CCCCCC}

% ── SECTION TITLES ──────────────────────────────────────────
\titleformat{\section}[block]
  {\normalfont\bfseries\large\color{amber}}
  {\thesection.}{0.5em}{}
  [\vspace{1pt}\hrule\vspace{2pt}]
\titlespacing{\section}{0pt}{10pt}{4pt}

\titleformat{\subsection}[runin]
  {\normalfont\bfseries\normalsize\color{cyanc}}
  {\thesubsection.}{0.4em}{}[ — ]
\titlespacing{\subsection}{0pt}{6pt}{2pt}

% ── BOX ENVIRONMENTS ────────────────────────────────────────
\newmdenv[
  backgroundcolor=lightbg,
  linecolor=amber,
  linewidth=1pt,
  innerleftmargin=8pt,
  innerrightmargin=8pt,
  innertopmargin=6pt,
  innerbottommargin=6pt,
  skipabove=5pt,
  skipbelow=5pt
]{amberbox}

\newmdenv[
  backgroundcolor=lightbg,
  linecolor=cyanc,
  linewidth=1pt,
  innerleftmargin=8pt,
  innerrightmargin=8pt,
  innertopmargin=6pt,
  innerbottommargin=6pt,
  skipabove=5pt,
  skipbelow=5pt
]{cyanbox}

\newmdenv[
  backgroundcolor=lightbg,
  linecolor=greenc,
  linewidth=1pt,
  innerleftmargin=8pt,
  innerrightmargin=8pt,
  innertopmargin=6pt,
  innerbottommargin=6pt,
  skipabove=5pt,
  skipbelow=5pt
]{greenbox}

\newmdenv[
  backgroundcolor=lightbg,
  linecolor=purplec,
  linewidth=1pt,
  innerleftmargin=8pt,
  innerrightmargin=8pt,
  innertopmargin=6pt,
  innerbottommargin=6pt,
  skipabove=5pt,
  skipbelow=5pt
]{purplebox}

\newmdenv[
  backgroundcolor=lightbg,
  linecolor=redc,
  linewidth=1pt,
  innerleftmargin=8pt,
  innerrightmargin=8pt,
  innertopmargin=6pt,
  innerbottommargin=6pt,
  skipabove=5pt,
  skipbelow=5pt
]{redbox}

% ── LISTS ───────────────────────────────────────────────────
\setlist[itemize]{leftmargin=14pt, itemsep=2pt, parsep=0pt, topsep=3pt}
\setlist[enumerate]{leftmargin=16pt, itemsep=2pt, parsep=0pt, topsep=3pt}

\setlength{\parindent}{0pt}
\setlength{\parskip}{3pt}

% ── MATH SHORTCUTS ──────────────────────────────────────────
\newcommand{\Tr}{T_r}
\newcommand{\dt}{\dot{s}}
\newcommand{\ddt}{\ddot{s}}

% ────────────────────────────────────────────────────────────
\begin{document}

{\centering
  {\Large\bfseries Trajectory Planning --- Cheat Sheet}\\[3pt]
  {\normalsize Robotics 1 $\cdot$ Prof.\ De Luca $\cdot$ Open Book}\\[2pt]
  \hrule
  \vspace{6pt}
}

% ════════════════════════════════════════════════════════════
\section{Diagramma Decisionale}

Rispondi in ordine prima di iniziare a calcolare.

\begin{amberbox}
\textbf{Q1 --- Dove si pianifica il moto?}\\
Il moto \`e specificato nel \textbf{joint space}? \quad$\Rightarrow$\quad \textbf{Caso A}\\
Il moto \`e specificato nel \textbf{Cartesian space}? \quad$\Rightarrow$\quad Q2

\medskip
\textbf{Q2 --- Hai un path geometrico o solo i due estremi?}\\
Solo punti iniziale/finale? \quad$\Rightarrow$\quad \textbf{Caso A}\\
Path geometrico dato (linea, cerchio, \ldots)? \quad$\Rightarrow$\quad Q3

\medskip
\textbf{Q3 --- I vincoli di velocit\`a/accelerazione sono su giunti o cartesiani?}\\
Bounds cartesiani diretti su $\|\dot{p}\|$, $\|\ddot{p}\|$? \quad$\Rightarrow$\quad \textbf{Caso B}\\
Bounds sui giunti $V_i$, $A_i$? \quad$\Rightarrow$\quad \textbf{Caso C}
\end{amberbox}

% ════════════════════════════════════════════════════════════
\section{Caso A --- Joint Space PTP}

\begin{cyanbox}
\begin{enumerate}
  \item Calcola $\Delta q_i = q_{f,i} - q_{i,i}$ per ogni giunto.
  \item Scegli il profilo: trapezoidale / cubica / quintica (vedi Sezioni 5--6).
  \item \textbf{Coordinazione:} $T = \max_i\{T_{\min,i}\}$; il giunto pi\`u lento detta il tempo.
  \item Gli altri giunti vengono scalati: \quad $k_i = T / T_{\min,i}$,\\
        \hspace*{1cm} $V_i^{\text{new}} = V_i / k_i$, \quad $A_i^{\text{new}} = A_i / k_i^2$.
  \item Scrivi $q_i(t)$ esplicitamente per ogni giunto.
\end{enumerate}
\end{cyanbox}

% ════════════════════════════════════════════════════════════
\section{Caso B --- Path Cartesiano con Bounds Diretti}

\begin{greenbox}
\begin{enumerate}
  \item Parametrizza $p(s)$ con parametro $s$ (lunghezza d'arco se possibile). Calcola $L$.
  \item Calcola $dp/ds$ e $d^2p/ds^2$ in funzione di $s$.
  \item Ricava $V$ dai bounds cartesiani diretti su $\|\dot{p}\|$.
  \item \textbf{Path lineare:} $\|\dot{p}\| = |\dt|$, $\|\ddot{p}\| = |\ddt|$. Bounds diretti.
  \item \textbf{Path circolare} (raggio $R$): vedi Sezione 7.
  \item Controlla $L$ rispetto a $V^2/A$ per scegliere tra b-c-b e bang-bang.
  \item Calcola $T$. Se serve $q(t)$: applica la cinematica inversa + Jacobiano.
\end{enumerate}
\end{greenbox}

% ════════════════════════════════════════════════════════════
\section{Caso C --- Bounds dai Giunti, Path Cartesiano}

\begin{purplebox}
\begin{enumerate}
  \item Applica la IK nei punti chiave: inizio, fine, e punti di \textbf{inversione}.
  \item Per ogni giunto, calcola la distanza totale percorsa.\\
        \textbf{Attenzione alle inversioni:} somma i tratti separatamente.\\
        \hspace*{1cm} $d_i = |q_{m,i} - q_{i,i}| + |q_{f,i} - q_{m,i}|$
  \item Stima del tempo minimo (accelerazione infinita):\quad $T_{\min,i} = d_i / V_i$.
  \item Giunto limitante: $i^* = \arg\max_i\{T_{\min,i}\}$.
  \item Worst-case geometrico: trova il punto del path dove $\|j_{i^*}\|$ \`e minimo.
  \item Bound sulla velocit\`a: \quad $V = \|j_{i^*}\|_{\min} \cdot V_{i^*}$.
  \item Poi procedi come il \textbf{Caso B}.
\end{enumerate}

\smallskip
\textit{Esempio (Jul 2024, robot RP):} $\|j_1\| = q_2$, minimo nel punto del path pi\`u vicino
all'origine. \quad $V = q_{2,\min} \cdot V_1$.
\end{purplebox}

\newpage
% ════════════════════════════════════════════════════════════
\section{Formule Trapezoidale (Bang-Coast-Bang)}

\begin{amberbox}
\textbf{Bang-Coast-Bang (b-c-b):} \quad condizione: $L > V^2/A$
\[
T_r = \frac{V}{A} \qquad\qquad T = \frac{L}{V} + \frac{V}{A} = \frac{AL + V^2}{AV}
\]

\textbf{Bang-Bang (se $L \leq V^2/A$):}
\[
\bar{T} = \sqrt{\frac{4L}{A}} \qquad\qquad \bar{V} = \frac{A\,\bar{T}}{2} \leq V_{\max}
\]

\textbf{Profilo $\sigma(t)$ esplicito (b-c-b):}
\[
\sigma(t) =
\begin{cases}
\dfrac{A\,t^2}{2} & 0 \leq t \leq T_r \\[8pt]
V\!\left(t - \dfrac{T_r}{2}\right) & T_r \leq t \leq T - T_r \\[8pt]
L - \dfrac{A\,(T-t)^2}{2} & T - T_r \leq t \leq T
\end{cases}
\]
\end{amberbox}

% ════════════════════════════════════════════════════════════
\section{Riscalamento Temporale}

\begin{cyanbox}
\textbf{Riscalamento uniforme} (fattore $k = \bar{T}/T \geq 1$):
\[
T_r^{\text{new}} = k\,T_r \qquad V^{\text{new}} = \frac{V}{k} \qquad A^{\text{new}} = \frac{A}{k^2}
\]

\textbf{Riscalamento con $T_r$ fisso} (esame Jan 2026):
\[
\bar{V} = \frac{L}{\bar{T} - T_r} \qquad \bar{A} = \frac{\bar{V}}{T_r}
\]

\textit{I due metodi danno risultati diversi. Leggi con attenzione il testo dell'esercizio.}
\end{cyanbox}

% ════════════════════════════════════════════════════════════
\section{Path Circolare --- Accelerazione Centripeta}

\begin{greenbox}
Per un path circolare di raggio $R$, parametrizzato ad arco $s$:
\[
\|\dot{p}\| = |\dt| \qquad\qquad
\|\ddot{p}\| = \sqrt{\ddt^2 + \frac{\dt^4}{R^2}} \leq a_{\max}
\]

In fase di \textbf{cruise} ($\ddt = 0$): \quad $\|\ddot{p}\| = \dfrac{V^2}{R} \leq a_{\max}$

In fase di \textbf{accelerazione} ($\dt = 0$): \quad $\|\ddot{p}\| = |\ddt| \leq a_{\max}$

\medskip
\textbf{Bounds ottimali:}
\[
V = \min\!\left\{v_{\max},\; \sqrt{R\,a_{\max}}\right\}
\qquad
A = \sqrt{a_{\max}^2 - \frac{V^4}{R^2}}
\]

\textit{Se $V = \sqrt{R\,a_{\max}}$ allora $A = 0$: non esiste soluzione in forma chiusa, \`e bang-bang.}
\end{greenbox}

% ════════════════════════════════════════════════════════════
\newpage
\section{Profili Polinomiali}

\subsection{Cubica (4 condizioni)}

Impone $q_i,\, q_f,\, \dot{q}(0) = v_i,\, \dot{q}(T) = v_f$.
\[
q(t) = a_0 + a_1 t + a_2 t^2 + a_3 t^3
\]
\[
a_0 = q_i \quad a_1 = v_i \quad
a_2 = \frac{3\Delta q - (2v_i + v_f)T}{T^2} \quad
a_3 = \frac{-2\Delta q + (v_i + v_f)T}{T^3}
\]

\textbf{Timing law rest-to-rest} ($v_i = v_f = 0$), in $\tau = t/T \in [0,1]$:
\[
s(\tau) = 3\tau^2 - 2\tau^3 \qquad \dot{s}(\tau) = \frac{6\tau(1-\tau)}{T}
\]
\[
\dot{s}_{\max} \text{ a } \tau=0.5:\quad
\begin{cases}
\dfrac{3}{2T} & \text{se } s\in[0,1] \\[6pt]
\dfrac{3L}{2T} & \text{se } s\in[0,L],\; s(\tau)=L(3\tau^2-2\tau^3)
\end{cases}
\]

\subsection{Quintica (6 condizioni)}

Aggiunge $\ddot{q}(0) = \ddot{q}(T) = 0$ alla cubica, in $\tau = t/T$:
\[
q(\tau) = q_i + \Delta q\,(10\tau^3 - 15\tau^4 + 6\tau^5)
\]
\[
|\dot{q}|_{\max} = \frac{30\,|\Delta q|}{16\,T} \quad \text{a } t = T/2
\qquad
|\ddot{q}|_{\max} = \frac{60\,|\Delta q|}{T^2}\cdot 0.0962 \quad \text{a } t \approx 0.21\,T
\]

Tempo minimo (coordinazione), con $|k| = 0.0962$:
\[
T^* = \max_i\!\left\{\,\frac{30\,|\Delta q_i|}{16\,V_{\max,i}},\quad
\sqrt{\frac{60\,|k|\,|\Delta q_i|}{A_{\max,i}}}\,\right\}
\]

% ════════════════════════════════════════════════════════════
\section{Velocit\`a nei Punti Chiave}

\begin{amberbox}
\begin{tabular}{@{}p{3.5cm}p{10cm}@{}}
$t = 0,\; T$ & $\dot{p} = 0 \;\Rightarrow\; \dot{q} = 0$ \quad (condizione rest-to-rest)\\[4pt]
$t = T/2$ & $\dt = V$;\quad $\dot{p}_m = V\,\dfrac{p_f - p_i}{L}$;\quad $\dot{q}_m = J^{-1}(q_m)\,\dot{p}_m$\\[6pt]
$t = T_r$ & $\dt = V$, $\ddt = 0$: solo accelerazione centripeta $V^2/R$ (path circolare)\\[4pt]
$t = 0^+$ & $\dt = 0$, $\ddt = A$: solo accelerazione tangenziale $A$\\
\end{tabular}
\end{amberbox}

\subsection{Emergency Stop da $t = T_0$}

Il robot inizia a decelerare massimamente da $T_0$:
\[
\dot{q}_i(t) = \dot{q}_i(T_0) - \operatorname{sign}\!\bigl(\dot{q}_i(T_0)\bigr)\cdot A_{\max,i}\,(t - T_0)
\]
\[
q_i(t) = q_i(T_0) + \dot{q}_i(T_0)\,(t-T_0) - \tfrac{1}{2}\,\operatorname{sign}(\dot{q}_i)\cdot A_{\max,i}\,(t-T_0)^2
\]
\[
T_{s,i} = T_0 + \frac{|\dot{q}_i(T_0)|}{A_{\max,i}} \qquad T_s = \max_i\{T_{s,i}\}
\]

% ════════════════════════════════════════════════════════════
\section{Trappole Comuni}

\begin{redbox}
\begin{itemize}
  \item \textbf{Giunto che inverte il moto:} $q_i = q_f$ non significa giunto fermo.
        La distanza percorsa \`e $|q_m - q_i| + |q_f - q_m| \neq 0$.

  \item \textbf{Accelerazione centripeta:} per path circolari $\|\ddot{p}\| \neq |\ddt|$.
        Il termine centripeto $\dt^4/R^2$ non \`e trascurabile ad alta velocit\`a.

  \item \textbf{Bang-bang dimenticato:} controlla \textit{sempre} la condizione $L$ vs $V^2/A$
        prima di usare la formula b-c-b.

  \item \textbf{Riscalamento uniforme vs $T_r$ fisso:} i due approcci danno risultati diversi.
        Non confonderli.

  \item \textbf{Quintica:} $T^*$ dipende da \textit{entrambi} i bounds $V$ e $A$.
        Non basta imporre solo il vincolo di velocit\`a.

  \item \textbf{Worst-case geometrico (robot RP):} $\|j_1\| = q_2$ \`e minimo nel punto
        del path pi\`u vicino all'origine, non necessariamente a met\`a arco.
\end{itemize}
\end{redbox}

% ════════════════════════════════════════════════════════════
\newpage
\section{Pattern dagli Esami (2024--2026)}

\begin{center}
\renewcommand{\arraystretch}{1.4}
\begin{tabular}{@{}p{2.6cm}p{11.5cm}@{}}
\toprule
\textbf{Esame} & \textbf{Pattern chiave}\\
\midrule
Jan 2026 & Trapezoidale con riscalamento a $T_r$ fisso. $\bar{V} = L/(\bar{T}-T_r)$, $\bar{A} = \bar{V}/T_r$.\\
Feb 2026 & Path circolare, scelta b-c-b vs b-b. $V = \min\{v_{\max}, \sqrt{Ra_{\max}}\}$, check $L$ vs $V^2/A$.\\
Jul 2025 & Path circolare bang-bang. $\|\ddot{p}\|_{\max} = \sqrt{A^2 + (V^2/R)^2}$ al punto medio.\\
Jun 2025 & Quintica coordinata + emergency stop. $T^* = \max\{T_{v,i}, T_{a,i}\}$ su tutti i giunti.\\
Jul 2024 & Path lineare robot RP, bounds da giunti. $V = q_{2,\min}\cdot V_1$. Giunto 2 inverte.\\
Sep 2024 & Coordinazione 2 giunti trapezoidali. $T = \max\{T_{\min,2}, T_{\min,3}\}$, scala il pi\`u veloce.\\
Feb 2025 & Trapezoidale con $v_i, v_f \neq 0$. Profilo b-c-b generalizzato; condizione cruise diversa.\\
Jun 2024 & Path in joint space $s\in[0,1]$. $\dot{q}(t) = q'(s)\,\dot{s}(t)$, $\dot{p}(t) = J(q)\,\dot{q}(t)$.\\
\bottomrule
\end{tabular}
\end{center}

\vspace{6pt}
\hrule
\vspace{4pt}
{\centering\small Robotics 1 $\cdot$ Prof.\ De Luca $\cdot$ Esami 2024--2026\par}

\end{document}
